\documentclass[11pt]{article}
\usepackage[a4paper,margin=25mm]{geometry}
\usepackage{amsmath,amssymb,amsthm,mathtools,booktabs,hyperref,xurl}
\setlength{\emergencystretch}{2em}
\newtheorem{theorem}{Theorem}
\newtheorem{lemma}[theorem]{Lemma}
\newtheorem{proposition}[theorem]{Proposition}
\newcommand{\Z}{\mathbb Z}\newcommand{\Q}{\mathbb Q}
\title{Signed power recovery and an absorbed Pell bound\\save six quadratic coordinates}
\author{\(\mathbb Z\mathbb Q\) Diophantine Project}
\date{23 November 2026}
\begin{document}
\maketitle

\begin{abstract}
The short Pell replacement for the Jones power introduces two strict
inequalities. Both separate positivity wrappers can be removed.
First, the retained coding slack and a congruence already force a large
absolute value of the proposed power; a signed estimate proves the
required prefix growth even when that power is initially negative.
The final Pell congruences then recover its positive value. Second,
the auxiliary index bound can be absorbed into an existing integral
small-Pell margin. These changes save six variables and two quadratic
residuals. Reusing an existing fourth-power register saves one further
variable and row in the ordinary-\(abc\) interface. The complete
general and ordinary-\(abc\) tables become respectively \(973/827\)
and \(697/592\). All integer interfaces and the cited coding theorem
are retained.
\end{abstract}

\section{The two changes and their scope}

Use the source and the two complete thirteen-input interfaces in
\cite{Power}. Set \(L=5^{60}\), \(T_b=b^5-1\), and retain
\[
 A_0=\frac{b^9+8b^5+b-8}{2},\qquad A=A_0+5T_bs,
 \qquad m_*=2bA-b^2-1.
\]
The auxiliary Pell certificate has integer inputs \(\delta,\alpha_*\)
and equations
\[
 C_*=L+\delta(A-1),\qquad
 D_*^2-(A^2-1)C_*^2=1,\qquad
 q=D_*-(A-b)C_*-\alpha_*m_*.
\]
Retain the cofinite instance range \(x>32z\), where all coding
parameters \(x,y,u,z\) are positive integers. Delete the separate
inequality \(q>b^5\). In Section~\ref{sec:absorbed-bound}, we also
delete the separate inequality \(C_*^2<c^2\) while strengthening
the existing small-Pell margin. Each deleted inequality had three
private positivity coordinates and one quadratic residual.

The first argument below temporarily retains \(C_*^2<c^2\).
The second argument restores that bound from the strengthened margin
before it is needed. This order makes the composition noncircular.
The claimed equivalence concerns the coding-prefix and instance
projection; private positivity representatives can be chosen anew.

\section{The proposed power has a large absolute value}

The preliminary integral reconstruction in \cite{Power} does not
require \(q>b^5\). It gives
\[
 A,C_*,D_*,q\in\Z,\qquad |A|\ge2,\qquad
 q\equiv\pm1\pmod{T_b}.
\]
Indeed \(A\equiv1\pmod5\), the congruence for \(C_*\) forces its
Pell index to be divisible by five, and the four signed Pell branches
reduce to \(\pm b^{\pm\nu}\) modulo \(T_b\).
The retained equation \(\lambda T_b=q^4-1\) therefore makes
\(\lambda\) integral even before the sign or exponent of \(q\)
is known.

Put \(Q=|q|\), and use the positive integral prefix inputs \(g,t,m\).
As in \cite{Power}, write
\[
 e=y+m(b^5-2z),\qquad
 \ell=u+t(b^5-2z),\qquad c_0=1+xb^5+g.
\]
The retained guards imply \(b\ge2\), \(e,\ell>0\), and
\begin{equation}\label{eq:slack}
                    e\ell g^2<(b-xy)q^2.
\end{equation}

\begin{lemma}\label{lem:large-Q}
Every proposed solution satisfies
\[
 b>xy\ge x>32z,\quad Q\ge b^5-2\ge b^4,\quad \lambda>0,
 \qquad e,\ell<\frac{Q^2}{b^3},\quad g<\frac{Q}{b^3},
 \qquad 0<c_0<bQ.
\]
\end{lemma}
\begin{proof}
The left side of \eqref{eq:slack} is positive, so \(q\ne0\) and
\(b>xy\). In particular \(2z<b\) and \(x\le b-1\). Hence
\[
             e,\ell\ge b^5-b+2\ge b^4.
\]
Since \(g\ge1\), equation~\eqref{eq:slack} gives \(Q^2>b^7\).
Thus \(Q>1\). The congruence \(q\equiv\pm1\pmod{T_b}\)
then implies \(Q\ge T_b-1=b^5-2\), and this is at least
\(b^4\) for \(b\ge2\). Consequently \(\lambda>0\).
Dividing \eqref{eq:slack} by the lower bounds for the other
positive factors gives \(e,\ell<Q^2/b^3\) and
\(g<Q/b^{7/2}<Q/b^3\).

As \(b^5\le Q+2\), we have
\[
 c_0<1+(b-1)b^5+Q/b^3
       \le(b-1)Q+(2b-1)+Q/b^3<bQ.
\]
The last step follows from
\(Q(1-b^{-3})\ge b^4-b\ge7b>2b-1\).
\end{proof}

\section{Signed prefix growth}

Retain the exact source expressions
\begin{align*}
 S={}&g+eq^3+\ell q^5+
   \left[2ec_0^4+\lambda\{b^5(1+q^4)-2zc_0^4\}\right]q^4,\\
 T={}&q^3-(b-1)\ell+(b^5-2z)\lambda q^3+(b^5-2)q^5,\\
 n={}&q^{16}=Q^{16},\qquad M=n^2=Q^{32},\qquad
 r=S(M-n)+T(M-1).
\end{align*}
The expression \(T\) need not be positive at this stage.

\begin{lemma}\label{lem:prefix}
Every proposed solution satisfies
\[
 S>14Q^7,\qquad |T|<2Q^7,\qquad
 r\ge M\ge b^{128}>2^{640}>L.
\]
\end{lemma}
\begin{proof}
Put \(E=b^5(1+Q^4)-2zc_0^4\). Since \(c_0<bQ\) and
\(2z<b/16\),
\[
 E>b^5+b^4Q^4(b-2z)>\frac{15}{16}b^5Q^4.
\]
Also
\(\lambda=(Q^4-1)/(b^5-1)>Q^4/(2b^5)\), since \(Q\ge2\).
Allowing either sign for the two odd powers in \(S\) gives
\[
 \begin{split}
 S&>\frac{15}{32}Q^{12}-eQ^3-\ell Q^5\\
  &>\frac{15}{32}Q^{12}-\frac{Q^5+Q^7}{b^3}
   \ge\left(15-\frac14\right)Q^7>14Q^7.
 \end{split}
\]
Here \(Q^5\ge32\), \(Q^5\le Q^7\), and \(b^3\ge8\).

For \(\theta=b^5-2z\), one has \(0<\theta<T_b\).
Therefore
\[
 \begin{split}
 |T|&\le Q^3+(b-1)\ell+\theta\lambda Q^3+(b^5-2)Q^5\\
     &<Q^7+\frac{Q^2}{b^2}+Q^6<2Q^7.
 \end{split}
\]
The middle inequality uses \(\theta\lambda<Q^4-1\) and
\(b^5-2\le Q\). The last follows already from \(Q,b\ge2\).

In particular \(S>2|T|+2\). As \(n\ge2\),
\[
 \begin{split}
 r&>(2|T|+2)(M-n)-|T|(M-1)\\
  &=|T|(n-1)^2+2(M-n)\ge M.
 \end{split}
\]
Finally \(Q\ge b^4\) gives \(M\ge b^{128}\), and
\(b>32z\ge32\) gives \(b^{128}>2^{640}>L\).
All quantities required to be integral in this reconstruction are
integral before any Pell index is identified.
\end{proof}

The retained signed binomial and half-trace arguments now apply:
they need integral \(r,M\), \(r\ge1\), \(M\ge16\), and
\(2\le b\le2M\), rather than a preliminary sign for \(q\).
They give the canonical positive \(w,s,c\),
\(bw=2^{2r+1}\), and the same growth estimates used in
\cite{Power}. Set \(U=wM\). In particular
\[
 s>U^{r-2},\quad a<s^2,\quad A-1>L+6r,\quad m_*>s,
 \quad C_*=Y_L(A)>0.
\]
The proof of the last conclusion uses the retained
\(C_*^2<c^2\); Section~\ref{sec:absorbed-bound} supplies it
from the final system's strengthened margin. Every use of
\(b<q\le U^{1/32}\) in the power proof is replaced by
\(b<Q=M^{1/32}\le U^{1/32}\).

\begin{theorem}[Exact recovery without a sign guard]\label{thm:signed-q}
The source with the separate inequality \(q>b^5\) deleted
still forces \(q=b^L\). It has the same coding-prefix and
instance projection as the predecessor.
\end{theorem}
\begin{proof}
Set \(K=b^L>1\). The two remaining trace signs in the auxiliary
Pell norm give, exactly as in \cite{Power},
\[
 D_*=X_L(A):\ q\equiv K\pmod{m_*},\qquad
 D_*=-X_L(A):\ qK\equiv-1\pmod{m_*}.
\]
Since \(b<Q\le U^{1/32}\) and \(L<r\),
\[
 QK+1\le2QK\le2U^{(6r+1)/32}<U^{r-2}<s<m_*.
\]
The strict middle comparison follows from
\((6r+1)/32+1\le r-2\) for \(r\ge4\) and \(U>2\).
The negative trace would make \(m_*\) divide \(qK+1\).
This integer is nonzero because \(q\in\Z\) and \(K>1\),
but has absolute value at most \(QK+1<m_*\), a contradiction.
For the positive trace,
\[
                  |q-K|\le Q+K\le QK+1<m_*.
\]
Its congruence therefore forces \(q=K=b^L>0\).
This equality restores the deleted inequality and its positivity
representatives. Conversely an old witness plainly satisfies the
system after that inequality is deleted.
\end{proof}

\section{The old integer margin absorbs the new index bound}
\label{sec:absorbed-bound}

Write \(C_*=L+\delta(A-1)\) for the new Pell coordinate and
retain the original source notation
\[
 Y=sM,\qquad p=2wr^2s^2M^3,\qquad
 \eta=k^2-4(c-kY)^2.
\]
Delete the separate inequality \(C_*^2<c^2\), and replace the
original requirement \(\eta>0\) by
\begin{equation}\label{eq:absorbed-margin}
                  \eta-C_*^2>0.
\end{equation}
This replacement needs no new coordinate: the square of the existing
coordinate \(C_*\) is already a quadratic monomial.

\begin{proposition}[Soundness of the absorbed bound]
Every rational solution of the modified source still satisfies the
old small-Pell strict inequality and the deleted index bound.
\end{proposition}
\begin{proof}
Equation~\eqref{eq:absorbed-margin} implies \(\eta>0\). Hence every
argument in the retained signed source that uses this inequality
remains valid. In particular, the half-trace square still excludes
\(s=0\) in the same noncircular order; this exclusion uses only
the original \(\eta>0\) consequence. Once \(s\ne0\), its
integrality and \(M\ge16\) give \(|Y|\ge16\). Since \(k\ne0\),
the old inequality then implies
\[
 \left|\frac{c}{k}-Y\right|<\frac12,
 \qquad |c|>|k|\left(|Y|-\frac12\right)>|k|.
\]
The strengthened margin also gives
\[
                 C_*^2<k^2-4(c-kY)^2\le k^2<c^2.
\]
Thus the deleted bound is available before the auxiliary Pell index
is identified. There is no dependence on knowing \(q=b^L\).
\end{proof}

\begin{proposition}[Completeness of the absorbed bound]
Every corrected positive source witness in the retained cofinite
instance range has a witness satisfying
\eqref{eq:absorbed-margin}, with the canonical auxiliary power
coordinate \(C_*=Y_L(A)\). The positivity coordinates can be chosen
anew; all source and Pell coordinates are retained.
\end{proposition}
\begin{proof}
The established positive reconstruction gives
\[
 w,s,k>0,\qquad k=Y_{r+1}(p),\qquad
 r\ge M\ge b^{128}>L,
 \qquad 0<A<5b^5s.
\]
The weaker lower bound \(M\ge b^{128}\) is sufficient here; the
canonical equality \(q=b^L\) provides a stronger bound. Since
\(w,r\ge1\), we obtain
\[
 p=2wr^2s^2M^3\ge2s^2M^3>100b^{10}s^2>4A^2.
\]
The strict middle inequality follows from
\(M^3\ge b^{384}>50b^{10}\), already true for \(b\ge2\).
For a positive Pell parameter the elementary recurrence bounds give
\[
 \begin{split}
 C_*^2=Y_L(A)^2
   &<(2A)^{2L-2}<p^{L-1}\\
   &<p^r<Y_{r+1}(p)=k.
 \end{split}
\]
Here \(L=5^{60}>2\), \(p\ge2\), and \(r>L-1\), so every
displayed strict inequality has its stated range.

It remains to compare this bound with the available margin. Put
\(e_0=c-kY\in\Z\). The original inequality is strict, so
\(k-2|e_0|\) is a positive integer. Consequently
\[
 \eta=(k-2|e_0|)(k+2|e_0|)\ge k>C_*^2.
\]
Thus \eqref{eq:absorbed-margin} holds. Its value is a positive
integer, so the same three triangular-number coordinates used by the
original positivity interface represent it. No additional integer
predicate or positivity variable is required.
\end{proof}

\paragraph{Exact row change.}
The retained positivity residual has the form
\[
       2\eta-2-\sum_{j=1}^3 u_j(u_j+1)=0.
\]
Subtracting the single quadratic monomial \(2C_*^2\) gives the
new residual. Delete the three private coordinates of the old
\(C_*^2<c^2\) wrapper and its single residual. This saves precisely
three coordinates and one quadratic residual, independently of the
three-coordinate saving from the signed-\(q\) argument.


\section{Complete tables and exact verification}

The two changes leave the thirteen integer inputs, every contact
family, every reciprocal guard, and all arithmetic-interface
hypotheses unchanged. In each complete predecessor table, delete
the six coordinates
\[
 \texttt{power\_positive\_q\_1,2,3},\qquad
 \texttt{power\_positive\_bound\_1,2,3}
\]
and their two private residuals. In the retained
\texttt{jones1980.positive.eta} residual, add exactly the monomial
\(-2\,\texttt{power\_C}^2\).

For the ordinary-\(abc\) interface there is one further exact
substitution. The retained rows
\[
 \texttt{int\_b\_29}=b^2,\qquad
 \texttt{int\_b\_30}=\texttt{int\_b\_29}^{\,2}
\]
already compute \(b^4\). Replace \texttt{power\_b4} by
\texttt{int\_b\_30} in the \(b^5\) and \(A\) definition rows,
and delete \texttt{power\_b4} and its duplicate defining row.
These are its only three occurrences. The substitution leaves every
residual quadratic and is exactly reversible by assigning
\texttt{power\_b4}=\texttt{int\_b\_30}. The general interface
stores projectively normalized powers and does not have this same
unscaled fourth-power register, so this additional saving is asserted
only for the ordinary-\(abc\) table. All other coefficients remain
unchanged.

\begin{center}
\begin{tabular}{lrrrr}
\toprule
Interface & Old variables & New variables & Old rows & New rows\\
\midrule
General & 979 & 973 & 829 & 827\\
Ordinary \(abc\) & 704 & 697 & 595 & 592\\
\bottomrule
\end{tabular}
\end{center}

The sum of squares of all residuals has degree exactly four and
vanishes over \(\Q\) exactly when every residual vanishes.
Exact degree follows from a retained residual with a nonzero quadratic
homogeneous part: its square cannot be cancelled in a real sum of squares.
The retained general family has its sufficiently large fixed
parameter as in \cite{Power}; no effective value is asserted.
The ordinary-\(abc\) family retains its stated arithmetic premise.
The source equivalence introduces no extra universality assumption
and retains the previously cited Jones coding theorem.

The builder and independent checker are
\begin{verbatim}
python undecidability/pell_signed_q_v98/build_signed_q.py
python undecidability/pell_signed_q_v98/verify_signed_q.py
\end{verbatim}
The checker imports no builder. It checks the exact deleted rows,
every occurrence of the six removed coordinates, the strengthened
margin, the two ordinary-\(abc\) fourth-power substitutions, all
826 or 589 unchanged rows, full support, degree, interface metadata,
predecessor hashes, and output certificates. The quantified arithmetic
arguments are the proofs above.

The independent \texttt{verify\_signed\_q\_lemmas.py} checks three
symbolic identities, exact signed-congruence and prefix-bound fixtures,
the integer margin inequality, signed guard recovery, and the Pell
growth estimates used in completeness. Its reproducible receipt is
\texttt{signed\_q\_lemmas\_certificate.json}. These finite checks audit
the algebra and edge cases; they do not replace the quantified proofs.

\begin{thebibliography}{9}
\small\setlength{\itemsep}{2pt}
\bibitem{Power} \(\mathbb Z\mathbb Q\) Diophantine Project,
\emph{A short Pell replacement for the fixed Jones power}.
\par\nolinkurl{undecidability/pell_power_v97/pell_power_replacement.tex}.
\end{thebibliography}
\end{document}
